Tabular models are routinely explained by running a sampling-based Shapley estimator once and reporting the top few attributes. When is that single run reliable evidence for the explanation of that prediction, and can the deployer tell from the run? We measure, per instance and at the level of the reported top-$k$ set, how often a run reproduces the exact interventional Shapley top-$k$, keeping repeatability across reruns apart from accuracy against the exact value, and we audit in-run certificates on fresh reruns. On two public tasks ($d=16$ and $d=20$) and three model classes, the shipped KernelSHAP default returns the exact top-3 set in 97\% of runs on average on ACS Employment, with a 7--10\% tail of near-tie instances, but in only 83--88\% on German Credit. There the error is repeatable, is not cured by sixteen times the budget, and comes from the default's $\ell_1$ feature selection; switching it off at the same budget gives 97--99\%. Four of five in-run certification rules are calibrated, but they are conservative, and cost at least about forty times what a fixed corrected estimator needs for the same exact rate, which gives no per-run statement.
